% ===============================
% MAT221 -- Real Analysis
% Riemann Integral
% ===============================

\documentclass[12pt,a4paper]{article}

% ===============================
% Metadata
% ===============================
\newcommand*{\CourseCode}{MAT221}
\newcommand*{\CourseTitle}{Real Analysis}
\newcommand*{\Chapter}{Properties of the Riemann Integral}
\newcommand*{\Topics}{ 
    Algebraic Properties, 
    Order Properties, 
    Additivity Property and Uniform Convergence in Riemann Integral
}
\newcommand*{\DocDate}{August 25, 2026}

% ===============================
% Header
% ===============================
\input{header}

% ===============================
% Document
% ===============================
\begin{document}

\FrontPage
\Large

\begin{heading}
    Cauchy Criterion for Riemann Integrability
\end{heading}

\vspace{10pt}

\begin{definition}[Riemann integrable function]
    A function $f$ is said to be Riemann integrable on the interval $[a,b]$ if for every $\epsilon > 0$, there exists a partition $P$ of $[a,b]$ such that
    \[
        U(f,P) - L(f,P) < \epsilon,
    \]
    where $U(f,P)$ and $L(f,P)$ are the upper and lower sums of $f$ with respect to the partition $P$, respectively.
    
\end{definition}

\clearpage

\begin{heading}
    Algebraic Properties of the Riemann Integral
\end{heading}

\vspace{10pt}

\begin{theorem}[Scalar multiplication]
    Let $f$ be Riemann integrable functions on the interval $[a,b]$ and $k \in \mathbb{R}$. Then, $kf$ is Riemann integrable on $[a,b]$, and
    \[
        \int_a^b kf = k \int_a^b f.
    \]
\end{theorem}

\clearpage

\begin{theorem}[Sum of integrable functions]
    Let $f$ and $g$ be Riemann integrable functions on the interval $[a,b]$. Then, $f+g$ is Riemann integrable on $[a,b]$, and
    \[
        \int_a^b (f + g) = \int_a^b f + \int_a^b g.
    \]
\end{theorem}

\clearpage

\begin{theorem}[Other algebraic properties]
    Let $f$ and $g$ be Riemann integrable functions on the interval $[a,b]$. Then, the following properties hold:
    \begin{enumerate}
        \item $f-g$ is Riemann integrable on $[a,b]$, and
        \[
            \int_a^b (f - g) = \int_a^b f - \int_a^b g.
        \]
        \item $fg$ is Riemann integrable on $[a,b]$.
        \item If $f(x) \neq 0$ for all $x \in [a,b]$, then $\frac{1}{f}$ is Riemann integrable on $[a,b]$.
        \item  If $g(x) \neq 0$ for all $x \in [a,b]$, then $\frac{f}{g}$ is Riemann integrable on $[a,b]$.
    \end{enumerate} 
\end{theorem}
    

\clearpage

\begin{heading}
    Order Properties of the Riemann Integral
\end{heading}

\vspace{10pt}

\begin{theorem}[Non-negativity]
    Let $f$ be a Riemann integrable function on the interval $[a,b]$. If $f(x) \geq 0$ for all $x \in [a,b]$, then
    \[
        \int_a^b f \geq 0.
    \]
\end{theorem}

\clearpage

\begin{theorem}[Order preservation]
    Let $f$ and $g$ be Riemann integrable functions on the interval $[a,b]$. If $f(x) \leq g(x)$ for all $x \in [a,b]$, then
    \[
        \int_a^b f \leq \int_a^b g.
    \]
\end{theorem}

\clearpage

\begin{theorem}[Absolute value]
    Let $f$ be a Riemann integrable function on the interval $[a,b]$. Then, $|f|$ is Riemann integrable on $[a,b]$, and
    \[
        \left| \int_a^b f \right| \leq \int_a^b |f|.
    \]
\end{theorem}

\clearpage

\begin{heading}
    Additivity of the Riemann Integral
\end{heading}

\vspace{10pt}

\begin{theorem}[Additivity]
    Let $f$ be a Riemann integrable function on the interval $[a,b]$ and $c \in (a,b)$. Then,
    \[
        \int_a^b f = \int_a^c f + \int_c^b f.
    \]
\end{theorem}

\begin{heading}
    Reversal of the limits of integration
\end{heading}

\begin{definition}[Integral from $b$ to $a$]
    For a Riemann integrable function $f$ on the interval $[a,b]$, the integral of $f$ from $b$ to $a$ is defined as
    \[
        \int_b^a f = - \int_a^b f .
    \]
\end{definition}

\vspace{100pt}

\begin{definition}[Integral over a point]
    For a Riemann integrable function $f$ on the interval $[a,b]$, the integral of $f$ over a point is defined as
    \[
        \int_a^a f = 0.
    \]
\end{definition}


\clearpage

\begin{heading}
    Uniform Convergence and the Riemann Integral
\end{heading}

\vspace{10pt}

\begin{theorem}[Uniform convergence and the Riemann integral]
    Let $\{f_n\}$ be a sequence of Riemann integrable functions on the interval $[a,b]$ that converges uniformly to a function $f$. Then, $f$ is Riemann integrable on $[a,b]$, and
    \[
        \lim_{n \to \infty} \int_a^b f_n = \int_a^b f.
    \]
    In other words,
    \[
        \lim_{n \to \infty} \int_a^b f_n = \int_a^b \lim_{n \to \infty} f_n.
    \]
\end{theorem}








\EndPage

\end{document}